% =====================================================================
%  Introduction generale
% =====================================================================
%  Le guide precise : « pose le probleme a resoudre sans evocation de
%  resultats », puis « le contenu de chaque chapitre est annonce
%  brievement ».
% =====================================================================
\cleardoublepage
\phantomsection
\addcontentsline{toc}{chapter}{Introduction generale}
\chapter*{Introduction generale}
\markboth{Introduction generale}{}

Le transport routier de marchandises occupe une place structurante dans
l'economie tunisienne. Il relie les zones industrielles du littoral aux
bassins de production de l'interieur, dessert les ports et assure la
distribution vers les reseaux commerciaux. Les entreprises du secteur
exploitent des flottes de plusieurs dizaines de vehicules, mobilisent des
chauffeurs sur des trajets longs et repondent a des clients qui attendent
des delais fermes et une facturation exacte.

Cette activite repose sur un equilibre fragile. Un vehicule immobilise
par une panne non anticipee, un document administratif expire decouvert
au moment d'un controle, un devis etabli sans reference tarifaire fiable
ou une facture oubliee : chacun de ces incidents a un cout direct et
degrade la relation commerciale. Les entreprises de taille moyenne
gerent pourtant ces risques avec des outils qui n'ont pas ete concus
pour cela, principalement des tableurs et des echanges informels.

\vspace{0.4cm}
\noindent
\essentiel{La question a laquelle ce projet de fin d'etudes entend
repondre est la suivante : comment doter une entreprise de transport de
marchandises d'un systeme d'information unifie qui centralise
l'exploitation de sa flotte, industrialise sa maintenance, ouvre un canal
numerique a ses clients et mette l'aide a la decision au service de ses
choix quotidiens ?}
\vspace{0.4cm}

Repondre a cette question suppose de traiter simultanement plusieurs
domaines metier qui, dans les pratiques observees, sont geres
separement : les identites et les habilitations, le referentiel des
vehicules et des equipements, le cycle commercial de la demande a la
facture, la maintenance curative et preventive, la communication interne
et la relation client. L'enjeu n'est pas seulement de les informatiser un
a un, mais de les faire cooperer sans les rendre indissociables : une
evolution de la facturation ne doit pas immobiliser la gestion de la
flotte.

Le travail presente ici a consiste a concevoir et a developper
\sujetCourt, une plateforme construite selon une architecture
microservices. Elle reunit douze services autonomes, trois applications
destinees a des publics distincts et une couche d'aide a la decision
fondee sur l'apprentissage automatique et sur la recherche documentaire
augmentee. La demarche suivie est agile, structuree par le cadre Scrum
et decoupee en six sprints correspondant chacun a un domaine metier
coherent.

\section*{Organisation du rapport}
\addcontentsline{toc}{section}{Organisation du rapport}

Ce rapport suit la progression reelle du projet.

Le \essentiel{premier chapitre} situe le travail dans son contexte. Il
presente l'organisme d'accueil, formule la problematique, examine les
pratiques et les solutions existantes, expose leurs limites, decrit la
solution proposee et justifie le choix de la methodologie de travail.

Le \essentiel{deuxieme chapitre} conduit l'analyse globale. Il identifie
les acteurs du systeme, recense les besoins fonctionnels et non
fonctionnels, les modelise par un diagramme de cas d'utilisation global,
puis presente le pilotage du projet : composition de l'equipe, carnet de
produit et decoupage en sprints.

Le \essentiel{troisieme chapitre} decrit l'environnement de travail,
materiel et logiciel, puis l'architecture retenue. Il en donne les deux
vues complementaires, logique et physique, et explique les trois choix
structurants qui la caracterisent.

Les \essentiel{chapitres quatre a neuf} rendent compte des six sprints.
Chacun suit la meme trame : planification et carnet de sprint,
elaboration technique appuyee sur les modeles UML, realisation et tests
de validation.

La \essentiel{conclusion generale} dresse le bilan du travail accompli,
revient sur les difficultes rencontrees, evalue l'atteinte des objectifs
et ouvre sur les perspectives d'evolution de la plateforme.
